\documentclass[11pt]{amsart}
\usepackage[margin=1.1in]{geometry}
\usepackage{lmodern}
\usepackage{amsmath,amssymb,amsthm,mathtools}
\usepackage[mathscr]{eucal}
\usepackage{booktabs,enumitem,xcolor,array}
\usepackage[colorlinks=true,linkcolor=blue!50!black,citecolor=green!40!black,urlcolor=blue!60!black]{hyperref}
\hypersetup{
  pdftitle={Extra prime atoms in positive solutions of the explicit formula: theta-group class sums and an obstruction},
  pdfauthor={Nic Johns}
}
\setlength{\emergencystretch}{2.5em}

% ---------------------------------------------------------------- macros
\newcommand{\RR}{\mathbb R}
\newcommand{\ZZ}{\mathbb Z}
\newcommand{\CC}{\mathbb C}
\newcommand{\NN}{\mathbb N}
\newcommand{\QQ}{\mathbb Q}
\newcommand{\dd}{\,\mathrm d}
\newcommand{\eps}{\varepsilon}
\renewcommand{\Re}{\operatorname{Re}}
\renewcommand{\Im}{\operatorname{Im}}
\newcommand{\sgn}{\operatorname{sgn}}
\newcommand{\one}{\mathbf 1}
\newcommand{\hF}{\widehat F}
\newcommand{\xin}[1]{\xi_{#1}}
\newcommand{\BRS}{\mathcal M}
\newcommand{\xiPP}{\xi_{\mathrm{PP}}}
\newcommand{\PP}{\mathrm{PP}}
\newcommand{\GR}{\Gamma_{\mathbb R}}
\newcommand{\Zz}{Z_\zeta}
\newcommand{\muz}{\mu_\zeta}
\newcommand{\nuz}{\nu_\zeta}
\newcommand{\pz}{p_\zeta}
\newcommand{\Oinf}{\Omega_\infty}
\newcommand{\Arch}{\mathcal A}
\newcommand{\Tc}{\mathcal T}
\newcommand{\Spec}{\mathcal R}
\newcommand{\Kad}{\mathcal K}
\newcommand{\uqpmat}[4]{\left(\begin{smallmatrix}#1&#2\\#3&#4\end{smallmatrix}\right)}

\newtheorem{theorem}{Theorem}[section]
\newtheorem{proposition}[theorem]{Proposition}
\newtheorem{lemma}[theorem]{Lemma}
\newtheorem{corollary}[theorem]{Corollary}
\newtheorem{fact}[theorem]{Fact}
\theoremstyle{definition}
\newtheorem{definition}[theorem]{Definition}
\theoremstyle{remark}
\newtheorem{remark}[theorem]{Remark}
\newtheorem{observation}[theorem]{Numerical observation}
\numberwithin{equation}{section}

\title[Extra prime atoms]{Extra prime atoms in positive solutions\\ of the explicit formula:
theta-group class sums and an obstruction}
\author{Nic Johns}
\date{October 2026}
\subjclass[2020]{Primary 11M26; Secondary 11F27, 42A38}
\keywords{Weil explicit formula, Fourier interpolation, theta group, modular integrals}

\begin{document}

\begin{abstract}
A companion paper studies positive solutions of Weil's explicit formula with the archimedean data of the Riemann zeta
function, and proves that a solution whose zero measure lives on the zeros of $\zeta$ plus finitely many extra points, and
whose prime measure lives on the interpolation nodes $\log m/4\pi$ of Bondarenko, Radchenko and Seip, is $\zeta$'s own. By a
theorem of Astra in the companion paper, a solution whose zero measure lives on the zeros of $\zeta$, possibly with an atom at the
origin, is $\zeta$'s own whatever its prime measure, so extra prime mass matters only together with extra zeros. This note treats extra atoms of the
prime measure off the nodes. Through theta-group class sums we show that an extra atom at a position $x$ with $x^2$ rational, if
it is the only one of this type, is excluded whenever an explicit Gauss-sum invariant $\kappa_{3,2}(x)$ is positive, also in the
presence of finitely many extra zeros, and that positions with $x^2$ irrational are invisible to every argument at leading order
in the class sums. Our statements for general non-residue classes and the nowhere-density of the possible positions assume no extra
zeros; they remain statements about theta class sums. On the zero side,
the same class sums give a positive-spanning property of the basis functions attached to the class $2\bmod3$.
\end{abstract}

\maketitle

\noindent\textbf{Status.} This note accompanies the companion paper \cite{PaperI} and has not been peer-reviewed. It records
partial results: every numbered statement in Sections~\ref{sec:classsums}, \ref{sec:signpatterns} and~\ref{app:proofs} is proved. The
question it addresses, uniqueness when the prime measure has extra atoms, is settled in \cite{PaperI} when there are no extra zeros
(\cite[Theorem~3.6]{PaperI}, due to Astra); the results here matter only when there are extra zeros, and that case remains open.
Countably many extra zeros, with summable weights, are also handled in \cite{PaperI} (Section~\ref{sec:countable}). The two numerical observations
(Section~\ref{sec:observations}) are not certified and are not used in any proof. The results of \cite{PaperI} that are used are
recalled in Section~\ref{sec:setting}, as Facts.

\setcounter{tocdepth}{1}
\tableofcontents


%======================================================================
\section{Introduction}\label{sec:intro}
%======================================================================

An \emph{admissible pair} is a pair of positive measures $(\mu,\nu)$, $\mu$ on the zero side and $\nu$ on the prime side
beyond $\xin2=\log2/2\pi$, that satisfies Weil's explicit formula with the archimedean data of $\zeta$ on a natural test class
(Section~\ref{sec:setting}). The companion paper \cite{PaperI} conjectures that every admissible pair, if one exists, is $\zeta$'s own
pair $\pz$, and proves this for pairs whose zero measure is carried by the real zeros of $\zeta$ together with finitely many extra
points and whose prime measure is carried by the set $\BRS=\{\log m/4\pi:m\ge1\}$ of the interpolation theorem of Bondarenko,
Radchenko and Seip \cite{BRS} (Fact~\ref{thm:Gfin}). Its proof tests the pair against the interpolation basis of \cite{BRS}, sums
the resulting node identities over the class $N\equiv2\pmod3$, which contains no squares, and concludes with a mean-value argument
for the resulting almost periodic function (Fact~\ref{lem:class2-landau}).

This note asks what happens when the prime measure has extra atoms off $\BRS$. An extra atom at $\eta$, with $x_\eta=e^{2\pi\eta}$,
enters the node identities through the coefficients $B_N(x_\eta)$ of a second modular integral $G_{x_\eta}$, and the natural tools
are its class sums over a quadratic non-residue class $b\bmod q$. We show:
\begin{enumerate}[label=\textup{(\roman*)}]
\item the class sums have a leading term $\kappa_{q,b}(x)v^{-1/2}$ as $v\to0^+$, with $\kappa_{q,b}(x)=0$ if $x^2\notin\QQ$ and $\kappa_{q,b}(x)$ an
explicit combination of normalised Gauss sums if $x^2\in\QQ$ (Proposition~\ref{prop:classsum});
\item rational-type extra atoms are constrained: if $\kappa_{q,b}(x_\eta)>0$ for some admissible class, a lone rational-type extra atom
at $\eta$ is excluded, whatever irrational-type extras accompany it (Theorem~\ref{thm:extra-primes});
\item the positions where $\kappa_{3,2}>0$ are dense (Lemma~\ref{lem:D}), so the positions that the class sums cannot exclude form a closed,
nowhere-dense set (Proposition~\ref{prop:C0});
\item irrational-type extras are invisible at leading order (Proposition~\ref{prop:invisible}): no argument at leading order in the class
sums can exclude them.
\end{enumerate}
Item (iv) is the obstruction. On the zero side, the same class sums show that the values of the BRS functions $U_m$,
$m\equiv2\pmod3$, at finitely many positive points off $\Zz$ positively span (Proposition~\ref{prop:signpattern}); this gives no
information on extra prime atoms. Numerical observations suggest that irrational-type extras are visible at intermediate scales
(Section~\ref{sec:observations}). The open problems are in Section~\ref{sec:open}.

The companion paper also contains two theorems contributed by Astra (OpenAI). The first, the zero-side
support theorem (Fact~\ref{fact:zerosupport}), shows that an admissible pair whose zero measure is carried by $\Zz\cup\{0\}$ is
$\zeta$'s pair, with no assumption on the prime measure. So every extra prime mass, atomic or not, at rational or irrational
positions, is excluded when there are no extra zeros, and the results of this note matter only when there are extra zeros: there
Theorem~\ref{thm:extra-primes}(b) and Proposition~\ref{prop:invisible} still apply, while Theorem~\ref{thm:extra-primes}(a) and
Proposition~\ref{prop:C0}(a), which assume no extra zeros, are now statements about the class sums only. Lemma~\ref{lem:P0}
anticipated the route of that theorem, whose proof carries it out in the limit, for all positions at once. The second theorem extends
Fact~\ref{thm:Gfin} to countably many extra zeros with summable weights and settles the extension proposed in
Section~\ref{sec:countable}.

%======================================================================
\section{Setting}\label{sec:setting}
%======================================================================

We recall the definitions and results of \cite{PaperI} that are used; statements quoted from \cite{PaperI} are called Facts.

\subsection{Admissible pairs}\label{ssec:pairs}
We write $\hF(\xi)=\int_\RR F(t)e^{-2\pi i\xi t}\dd t$ and $\xin n=\log n/(2\pi)$ for real $n>0$. The archimedean functional is
$\Arch(F):=F(\frac i2)+F(-\frac i2)+\int_\RR F\,\Oinf$, with $\Oinf(t)=\frac1{2\pi}(\Re\psi(\frac14+\frac{it}2)-\log\pi)$. The test class
$\Tc=\bigcup_{0<\delta<1/2}\Tc_\delta$ consists of the even functions $F$, real on $\RR$ and analytic in a neighbourhood of
$\{\lvert\Im z\rvert\le\frac12+\delta\}$ for some $\delta$, with $\sup(1+\lvert z\rvert)^2\lvert F(z)\rvert<\infty$ there.

\begin{fact}[Explicit formula on $\Tc$ {\cite{PaperI}}; {\cite{Wei52}}, {\cite[(1.1)]{BRS}}]\label{lem:explicit-formula}
For every $F\in\Tc$, $\Arch(F)=\sum_\rho F(t_\rho)+\frac1\pi\sum_{n\ge2}\Lambda(n)n^{-1/2}\hF(\xin n)$, where $t_\rho:=-i(\rho-\frac12)$ and
$\rho$ runs over the non-trivial zeros of $\zeta$ with multiplicity. No hypothesis on the zeros is made.
\end{fact}

$\Zz$ is the set of real ordinates $\gamma$ with $\zeta(\frac12+i\gamma)=0$, $m_\gamma$ the multiplicity,
$\muz=\sum_{\gamma\in\Zz}m_\gamma\delta_\gamma$, $\nuz=\sum_{n\ge2}\Lambda(n)n^{-1/2}\delta_{\xin n}$, and $\pz=(\muz,\nuz)$. An
\emph{admissible pair} is a pair $(\mu,\nu)$ of positive Radon measures, $\mu$ even on $\RR$ and $\nu$ on $[\xin2,\infty)$, such
that $\Arch(F)=\Spec_{\mu,\nu}(F):=\int F\dd\mu+\frac1\pi\int\hF\dd\nu$ for every $F\in\Tc$, with absolutely convergent integrals; $\Kad$
is the set of admissible pairs. We write $\BRS:=\{\log m/(4\pi):m\ge1\}$; since $\xin n=\log(n^2)/(4\pi)$, the prime-power points
$\xiPP=\{\xin n:n\in\PP\}$ lie in $\BRS$.

$\vartheta(\tau)=\sum_{n\in\ZZ}e^{\pi in^2\tau}$ is Jacobi's theta function and $\Gamma_\vartheta$ the theta group, generated by
$S=\uqpmat{0}{-1}{1}{0}$ and $T^2=\uqpmat1201$ and viewed, as in \cite{BRS}, in $\mathrm{PSL}_2(\RR)$, so that matrices are identified up to
sign and $S^2=1$; we write $\tau=\sigma+iv\in\mathbb H$. $\boldsymbol\mu$ is the M\"obius function,
$\one_\square(m)=1$ if $m$ is a perfect square and $0$ otherwise, and $\Lambda(\sqrt m):=0$ if $m$ is not a square.

\subsection{The interpolation basis}\label{ssec:basis}
For $s\in\CC$ let $\mathsf F(\tau,s)=\sum_{N\ge0}\alpha_N(s)e^{\pi iN\tau}$ be the unique $2$-periodic analytic function of
moderate growth on $\mathbb H$ satisfying
\begin{equation}\label{eq:BRS-FE}
\mathsf F(\tau,s)+(\tau/i)^{-1/2}\,\mathsf F(-1/\tau,s)=\mathfrak p_s(\tau):=(\tau/i)^{-s}+(\tau/i)^{s-1/2}
\end{equation}
(principal branches); this is $F^-_{1/2}(\tau,s)$ of \cite[Theorem~3.1]{BRS}. Its coefficients grow polynomially in $N$, are entire
in $s$, and are real on the line $\Re s=\frac14$. Let $b_N$ $(N\ge0)$ be the even Fourier eigenfunctions with eigenvalue $+1$ of
Radchenko and Viazovska \cite{RV19} (the functions $b_N^-$ of \cite[\S3]{BRS}, cf.\ \cite[\S3.4.1]{BRS}), so that
$b_N(\sqrt n)=\delta_{Nn}$ for $n\ge1$, and put $B_N(x):=\sum_{k\ge1}b_N(kx)$ for $x>0$. By \cite[Theorem~1.1]{BRS}, which assumes
neither RH nor simplicity of the zeros, there are even entire functions $U_m$ ($m\ge1$) and $V_{\rho,j}$ ($\rho$ a non-trivial zero
with $\Im\rho>0$, $0\le j<m(\rho)$) with $U_m^{(j)}(t_\rho)=0$, $\widehat{U_m}(\log m'/4\pi)=\delta_{mm'}$,
$V_{\rho,j}^{(j')}(t_{\rho'})=\delta_{(\rho,j),(\rho',j')}$ and $\widehat{V_{\rho,j}}(\log m/4\pi)=0$; their real and imaginary parts lie in
$\Tc$. These are the \emph{BRS functions}. With $\Xi(t):=\xi(\frac12+it)$, $\xi(s)=\frac12s(s-1)\GR(s)\zeta(s)$ and $\GR(s)=\pi^{-s/2}\Gamma(\frac s2)$,
explicitly \cite[\S\S4.2--4.3]{BRS}
\begin{equation}\label{eq:Um}
U_m(z)=\frac{m^{-1/4}}{2\pi}\,h_m\big(\tfrac12+iz\big),\qquad h_m(s)=\frac{\GR(s)\zeta(s)}2\sum_{d^2\mid m}\boldsymbol\mu(d)\,\alpha_{m/d^2}(s/2).
\end{equation}
and $\widehat{U_m}(\xi)=m^{-1/4}e^{\pi\xi}u_m(e^{2\pi\xi})$ with $u_m(x)=\sum_{d^2\mid m}\boldsymbol\mu(d)B_{m/d^2}(x)$. By the explicit formula,
applied to real and imaginary parts,
\begin{equation}\label{eq:W-BRS}
\Arch(U_m)=\frac1\pi\Lambda(\sqrt m)\,m^{-1/4},\qquad \Arch(V_{\rho,0})=2m(\rho),\qquad \Arch(V_{\rho,j})=0\quad(j\ge1),
\end{equation}
where $m(\rho)$ is the multiplicity of $\rho$ and $\Lambda(\sqrt m):=0$ if $m$ is not a square \cite{PaperI}. Indeed $U_m$ vanishes at every
$t_\rho$ and $\widehat{U_m}(\xin n)=\widehat{U_m}(\log(n^2)/4\pi)=\delta_{m,n^2}$; for $V_{\rho,j}$ the prime sum vanishes, and the zeros with
negative imaginary part contribute as their conjugates do, by evenness.

Fix a finite symmetric set $E\subset\RR\setminus\Zz$, and put $E_+:=E\cap[0,\infty)$ and $s_e:=\frac14+\frac{ie}2$. Let
$\mathcal X\subset[\xin2,\infty)\setminus\BRS$ be finite, and write $x_\eta:=e^{2\pi\eta}$ for $\eta\in\mathcal X$; so $x_\eta\ge2$
and $x_\eta^2\notin\NN$.

\begin{lemma}[Structure of pairs with extra atoms]\label{lem:structure}
Let $\mu$ be an even real measure carried by $\Zz\cup E$ and $\nu$ a real measure carried by $\BRS\cup\mathcal X$ such that
$\Arch(F)=\Spec_{\mu,\nu}(F)$ for $F=U_m$ $(m\ge1)$ and $F=V_{\rho,j}$. Write $a_e:=\mu(\{e\})$, $w_\gamma:=\mu(\{\gamma\})$,
$\beta_\eta:=\nu(\{\eta\})$, $\ell_m:=m^{1/4}\nu(\{\frac{\log m}{4\pi}\})$, $C_N:=\sum_{k^2\mid N}\ell_{N/k^2}$,
\[
\tilde\varrho_e:=\frac{a_e\,\Xi(e)}{\frac14+e^2}\ (e>0),\qquad \tilde\varrho_0:=2a_0\,\Xi(0),\qquad c_N(E):=\sum_{e\in E_+}\tilde\varrho_e\,\alpha_N(s_e)\in\RR .
\]
Then for every $N\ge1$
\begin{equation}\label{eq:CN}
C_N=c_N(E)+\tfrac12(\log N)\one_\square(N)-\sum_{\eta\in\mathcal X}\beta_\eta\sqrt{x_\eta}\,B_N(x_\eta).
\end{equation}
If $\mathcal X=\emptyset$, moreover: $c_1=c_2=c_3=0$ when $\nu$ is carried by $[\xin2,\infty)$; $w_\gamma=m_\gamma-\frac12\sum_{e\in E}a_eV_{\rho,0}(e)$
for $\gamma>0$ and $\rho=\frac12+i\gamma$; $\frac12\sum_ea_eV_{\rho,0}(e)=m(\rho)$ for every zero $\rho$ off the line with $\Im\rho>0$;
$\sum_ea_eV_{\rho,j}(e)=0$ for $1\le j<m(\rho)$; and $(\mu,\nu)=\pz$ if and only if $a=0$, if and only if $c(E)\equiv0$. If
$(\mu,\nu)\in\Kad$, then $C_N\ge0$ for all $N$; in particular $c_N(E)\ge0$ for every non-square $N$ when $\mathcal X=\emptyset$.
\end{lemma}

The proof is in Appendix~\ref{app:proofs-structure}. The case $\mathcal X=\emptyset$, with the statements that follow ``moreover'',
is \cite[Lemma~3.7]{PaperI}, which allows countable $E$.



\begin{fact}[Finitely many extra zeros; the case of finite $E$ of {\cite[Theorem~3.8]{PaperI}}]\label{thm:Gfin}
Let $E\subset\RR\setminus\Zz$ be finite and symmetric. Let $\mu$ be an even real measure carried by $\Zz\cup E$ and $\nu$ a real
measure carried by $\BRS$ such that $\Arch(F)=\int F\dd\mu+\frac1\pi\int\hF\dd\nu$ for every BRS function $F$ (in particular, if the
identity holds on $\Tc$ with absolutely convergent integrals). Assume $\nu(\{\frac{\log m}{4\pi}\})\ge0$ for every $m\equiv2\pmod3$, and
$\mu(\{0\})\ge0$ if $0\in E$. Then RH holds, $\mu=\muz$ and $\nu=\nuz$. In particular every admissible pair carried by
$(\Zz\cup E)\times\BRS$ is $\pz$ \textup{(}\cite[Theorem~B(b)]{PaperI}\textup{)}.
\end{fact}

In \cite[Theorem~3.8]{PaperI} this is proved, by a theorem of Astra, for countable symmetric $E\subset\RR\setminus\Zz$ with
$\sum_{e\in E,\,e>0}\lvert\mu(\{e\})\rvert\,\lvert\zeta(\frac12+ie)\rvert\,(1+e)^{-1/2}<\infty$, or more generally with partial sums of
$\lvert\mu(\{e\})\rvert\,\lvert\zeta(\frac12+ie)\rvert$ over $0<e\le T$ of order $o(T^{1/2})$.

\begin{fact}[Zero-side support; {\cite[Theorem~3.6]{PaperI}}, due to Astra]\label{fact:zerosupport}
If $(\mu,\nu)\in\Kad$ and $\mu$ is carried by $\Zz\cup\{0\}$, then RH holds and $(\mu,\nu)=\pz$. No assumption on $\nu$ is made.
\end{fact}

\subsection{The theta function near rational points}\label{ssec:theta}
A fraction $a/c$ in lowest terms ($c\ge1$) has one of the three \emph{types} (even, odd), (odd, even) and (odd, odd), according
to the parities of $a$ and $c$. For $a/c\in\QQ$ in lowest terms let $g(a/c):=c^{-1}\sum_{n\bmod c}e^{\pi ian^2/c}$ if $ac$ is even, and
$g(a/c):=0$ if $ac$ is odd.
We use $\vartheta(-1/\tau)=(\tau/i)^{1/2}\vartheta(\tau)$, $\vartheta(\tau+2)=\vartheta(\tau)$, and the fact that $\vartheta$ has no zeros in $\mathbb H$
(see \cite[\S3.2]{BRS}). Principal branches are used throughout; $(\tau/i)^\alpha>0$ for $\tau\in i(0,\infty)$.

\begin{lemma}[$\vartheta$ near real points]\label{lem:theta-local}
\begin{enumerate}[label=\textup{(\alph*)}]
\item Let $r=a/c\in\QQ$ in lowest terms, $c\ge1$. For $z\in\mathbb H$ with $Y:=\Im(-1/z)\ge1$,
$\vartheta(r+z)=g(r)(z/i)^{-1/2}+O_r(|z|^{-1/2}e^{-\pi Y/(4c^2)})$. In particular, if $z\to0$ with $|\Re z|\le M\Im z$, then
$\vartheta(r+z)=g(r)(z/i)^{-1/2}+O(|z|^{-1/2}e^{-c'/|z|})$ for some $c'=c'(r,M)>0$.
\item Let $r\in\RR\setminus\QQ$. If $\eps\to0^+$ and $|\delta|\le M\eps^2$, then $\vartheta(r+\delta+i\eps)=o(\eps^{-1/2})$.
\item For $r\in\QQ\setminus\{0\}$, $g(-1/r)=|r|^{-1/2}e^{-i\pi\sgn(r)/4}\,g(r)$.
\end{enumerate}
\end{lemma}

The proof is in Appendix~\ref{app:proofs-theta}.

\subsection{Cusp expansions}\label{ssec:cusp}\label{app:cusp}
Let $\chi:\Gamma_\vartheta\to\{\pm1\}$ be the character with $\chi(S)=-1$ and $\chi(T^2)=1$; it is well defined because $S$ and $T^2$ generate
$\Gamma_\vartheta$ with the single relation $S^2=1$ \cite[\S3.1]{BRS}. For $f:\mathbb H\to\CC$ and $\gamma\in\Gamma_\vartheta$ put
\[
(f|\gamma)(\tau):=\chi(\gamma)\,\frac{\vartheta(\tau)}{\vartheta(\gamma\tau)}\,f(\gamma\tau),
\]
a right action: $(f|\gamma_1)|\gamma_2=f|(\gamma_1\gamma_2)$. As $\vartheta(\tau)/\vartheta(-1/\tau)=(\tau/i)^{-1/2}$, the functional equations
\eqref{eq:BRS-FE} and \eqref{eq:Gx-FE} and periodicity read $f-f|S=\mathfrak p$ and $f|T^2=f$, for $(f,\mathfrak p)=(\mathsf F(\cdot,s),\mathfrak p_s)$ and
$(G_x,\mathfrak p_x)$. If $\gamma=g_m\cdots g_1$ with $g_k\in\{S,T^{\pm2}\}$ and $w_k:=g_k\cdots g_1$ ($w_0=1$), then telescoping
$f|w_{k-1}-f|w_k=(f-f|g_k)|w_{k-1}$ gives
\begin{equation}\label{eq:unroll}
f=f|\gamma+\sum_{k:\,g_k=S}\mathfrak p|w_{k-1}.
\end{equation}
For the cusp $\frac23$ we use $\gamma=ST^2ST^2S=\uqpmat{-2}{1}{3}{-2}$, which maps $\frac23$ to $\infty$, with $S$ in positions $k=1,3,5$ and
$w_2=T^2S$, $w_4=T^2ST^2S=\uqpmat{3}{-2}{2}{-1}$. For $\tau=\frac23+iv$ one finds $T^2S\tau=2-1/\tau\to\frac12$, $w_4\tau=9iv/(1+6iv)$ and
$\gamma\tau=-\frac23+\frac i{9v}$; moreover $\chi(\gamma)=-1$, $\chi(T^2S)=-1$, $\chi(w_4)=1$,
\begin{equation}\label{eq:jfactors}
\frac{\vartheta(\tau)}{\vartheta(T^2S\tau)}=(\tau/i)^{-1/2},\qquad
\frac{\vartheta(\tau)}{\vartheta(w_4\tau)}=(\tau/i)^{-1/2}\big((2-1/\tau)/i\big)^{-1/2},
\end{equation}
and both are analytic near $\tau=\frac23$, the second with value $(\frac32)^{1/2}e^{i\pi/4}\cdot\sqrt2\,e^{i\pi/4}=i\sqrt3$ there.

\begin{fact}[The class $2\bmod3$; {\cite[Lemma~B.2]{PaperI}}]\label{lem:cusp-class2}
Let $\Re s=\frac14$. As $v\to0^+$,
\begin{multline*}
\tfrac13\Big[\mathsf F(iv,s)+\omega\,\mathsf F(\tfrac23+iv,s)+\bar\omega\,\mathsf F(\tfrac43+iv,s)\Big]\\
=\tfrac13\big(1-3\cdot9^{-s}\big)v^{-s}+\tfrac13\big(1-3\cdot9^{s-1/2}\big)v^{s-1/2}+c(s)+O_s(v^{3/4}),
\end{multline*}
with a constant $c(s)$. Consequently, if $E_+\subset[0,\infty)$ is finite, $\tilde\varrho_e\in\RR$ and $c_N:=\sum_{e\in E_+}\tilde\varrho_e\alpha_N(s_e)$, then
\[
\Theta_E^{(2)}(v):=\sum_{N\equiv2\,(3)}c_Ne^{-\pi Nv}=\sum_{e\in E_+,\,e>0}\big(B_ev^{-s_e}+\overline{B_e}v^{-\overline{s_e}}\big)+B_0v^{-1/4}+C+O(v^{3/4}),
\]
with $B_e=\frac{\tilde\varrho_e}3(1-3\cdot9^{-s_e})$ for $e>0$, $B_0=\frac{2\tilde\varrho_0}3(1-\sqrt3)$ (and $B_0=0$ if $0\notin E_+$), and a real constant $C$. In
particular $\Theta_E^{(2)}(v)=O(v^{-1/4})$.
\end{fact}

The function $\Theta_E^{(2)}$ is the class sum of the node identities of Lemma~\ref{lem:structure} with $\mathcal X=\emptyset$ over the class
$2\bmod3$; in \cite{PaperI} its expansion and a mean-value argument give Fact~\ref{thm:Gfin}.

\begin{fact}[Mean-value step; the case of finite $E_+$ of {\cite[Lemma~B.4]{PaperI}}]\label{lem:class2-landau}
Let $E_+\subset[0,\infty)$ be finite, let $a_e\in\RR$ $(e\in E_+)$, put $\tilde\varrho_e:=a_e\Xi(e)/(\frac14+e^2)$ for $e>0$ and
$\tilde\varrho_0:=2a_0\Xi(0)$, and $c_N:=\sum_{e\in E_+}\tilde\varrho_e\alpha_N(s_e)$. If $a_0\ge0$ when $0\in E_+$, and $c_N\ge0$ for every
$N\equiv2\pmod 3$, then $\tilde\varrho_e=0$ for every $e\in E_+$.
\end{fact}

%======================================================================
\section{Class sums of an extra atom}\label{sec:classsums}\label{ssec:extraprimes}

We turn to admissible pairs whose prime measure has atoms off $\BRS$, keeping the notation of Lemma~\ref{lem:structure}. An
extra atom at $\eta$ enters the node identities \eqref{eq:CN} through the numbers $B_N(x_\eta)$, and these are the Fourier
coefficients of a second modular integral.

\begin{definition}[Class sums of an extra atom]\label{def:classsum}
For $x>0$ and $\tau\in\mathbb H$ put $G_x(\tau):=\sum_{N\ge0}B_N(x)e^{\pi iN\tau}$. This is the modular integral
$F^-_{1/2}(\tau,\phi_x)$ of \cite[Theorem~3.1]{BRS} attached to $\phi_x(\tau):=\frac12(\vartheta(x^2\tau)-1)$ (cf.\ the proof of
\cite[Theorem~1.1]{BRS}, where the argument of $\vartheta$ is printed as $zx$; the square is needed, since
$B_N(\sqrt m)=\#\{k\ge1:k^2m=N\}$ for integers $m\ge1$): it is $2$-periodic, of moderate growth, and
\begin{equation}\label{eq:Gx-FE}
G_x(\tau)+(\tau/i)^{-1/2}G_x(-1/\tau)=\mathfrak p_x(\tau):=\tfrac12\Big[\vartheta(x^2\tau)-1+x^{-1}\vartheta(\tau/x^2)-(\tau/i)^{-1/2}\Big].
\end{equation}
For an odd prime $q$, a quadratic non-residue $b\bmod q$ and $v>0$ let
\[
T^{(x)}_{q,b}(v):=\sum_{N\equiv b\,(q)}B_N(x)e^{-\pi Nv}=\frac1q\sum_{j\bmod q}e^{-2\pi ijb/q}\,G_x\big(\tfrac{2j}q+iv\big).
\]
\end{definition}

Recall that the number $g(r)$ of Section~\ref{ssec:theta} describes $\vartheta$ near the rational point $r$:
$\vartheta(r+z)=g(r)(z/i)^{-1/2}+O(|z|^{-1/2}e^{-c/|z|})$ as $z\to0$ non-tangentially in $\mathbb H$ (Lemma~\ref{lem:theta-local}).

\begin{proposition}[Leading term of non-residue class sums]\label{prop:classsum}
Let $x>0$, $q$ an odd prime and $b$ a quadratic non-residue mod $q$. As $v\to0^+$,
\[
T^{(x)}_{q,b}(v)=\kappa_{q,b}(x)\,v^{-1/2}+o(v^{-1/2}),
\]
where $\kappa_{q,b}(x)=0$ if $x^2\notin\QQ$, and for $x^2\in\QQ$, $\kappa_{q,b}(x)$ is an explicit finite linear combination of values of
$g$, the error being $O(1)$. For $(q,b)=(3,2)$, with $X:=x^2\in\QQ$ and $\omega:=e^{2\pi i/3}$,
\[
\kappa_{3,2}(x)=\tfrac23\Re(\omega K),\qquad K=\frac{g(\frac{2X}3)}{2x}+\frac{e^{i\pi/4}}{\sqrt6\,x}\Big[g\big(-\tfrac{3X}2\big)-g\big(\tfrac X2\big)\Big]-\frac{i\,g(-2X)}{2\sqrt3\,x}.
\]
For $X=m\in\NN$ this gives $\kappa_{3,2}(x)=1/(3x)$ if $m\equiv2\pmod3$ and $\kappa_{3,2}(x)=0$ otherwise.
\end{proposition}

The proof, in Appendix~\ref{app:proofs-classsum}, unrolls \eqref{eq:Gx-FE} along a reduction word at each cusp $2j/q$. The
contribution $(B_0(x)-\frac1{2x})\vartheta(\tau)$ that a single atom would make to every class appears at each cusp with the same
coefficient, because the multiplier of the reduction word is $-1$ at every cusp with odd denominator; it therefore cancels
exactly in every non-residue class, since $\sum_je^{-2\pi ijb/q}\vartheta(\frac{2j}q+iv)=q\sum_{n^2\equiv b\,(q)}e^{-\pi n^2v}=0$. What remains
samples $\vartheta$ near the points $x^2\rho$ and $\rho/x^2$ for finitely many rationals $\rho$. The case $X=m\in\NN$ is a check:
then $B_N(\sqrt m)=\sum_k\delta_{N,k^2m}$, so $G_x=\phi_x$ and $T^{(x)}_{3,2}(v)=\sum_{3\nmid k,\ k^2m\equiv2\,(3)}e^{-\pi k^2mv}$.

\begin{lemma}[Node-positive test functions suffice]\label{lem:P0}
Let $(\mu,\nu)\in\Kad$ be carried by $(\Zz\cup E)\times(\BRS\cup\mathcal X)$ as in Lemma~\ref{lem:structure}, and let $\eta_0\in\mathcal X$.
Suppose some $F\in\Tc$ (or the real part of an even complex $F$ whose real and imaginary parts lie in $\Tc$) vanishes at every
$t_\rho$ and on $E$, has $\hF\ge0$ on $(\BRS\cap[\xin2,\infty))\cup\mathcal X$, $\hF=0$ on $\xiPP$, and $\hF(\eta_0)>0$. Then $\beta_{\eta_0}=0$.
\end{lemma}

\begin{proof}
By the explicit formula (Fact~\ref{lem:explicit-formula}), $\Arch(F)=0$. On the spectral side $\int F\dd\mu=0$ and
$\frac1\pi\int\hF\dd\nu\ge\frac1\pi\beta_{\eta_0}\hF(\eta_0)\ge0$.
\end{proof}

No positivity of $F$ on $\RR$, and none of $\hF$ off $\BRS\cup\mathcal X$, is needed. Examples are $\Xi^2H$ and non-negative
combinations of BRS functions $U_m$ over non-square $m$. When $E\subseteq\{0\}$ the route of this lemma is carried out, in the limit and
for all positions at once, by the proof of Fact~\ref{fact:zerosupport}.

\begin{theorem}[Rational prime-side extras]\label{thm:extra-primes}
Let $(\mu,\nu)\in\Kad$ with $\mu$ carried by $\Zz\cup E$ and $\nu$ carried by $\BRS\cup\mathcal X$, where $E$ and $\mathcal X$ are finite as
in Lemma~\ref{lem:structure}. Let $\mathcal X_\QQ:=\{\eta\in\mathcal X:x_\eta^2\in\QQ\}$.
\begin{enumerate}[label=\textup{(\alph*)}]
\item If $E=\emptyset$, then for every odd prime $q$ and every quadratic non-residue $b\bmod q$,
$\sum_{\eta\in\mathcal X_\QQ}\beta_\eta\sqrt{x_\eta}\,\kappa_{q,b}(x_\eta)\le0$.
\item For arbitrary finite $E$, the inequality holds for $(q,b)=(3,2)$.
\end{enumerate}
Consequently, if some non-negative combination of the $\kappa_{q,b}$ admissible in \textup{(a)} or \textup{(b)} is positive at every
$x_\eta$, $\eta\in\mathcal X_\QQ$, then $\beta_\eta=0$ for all $\eta\in\mathcal X_\QQ$. In particular, if $\eta$ is the only rational-type extra atom and
$\kappa_{q,b}(x_\eta)>0$ for one admissible $(q,b)$, then $\beta_\eta=0$, whatever irrational-type extras accompany it. (With
several rational-type extras the inequalities of (a) and (b) are signed sums, and individual terms may have opposite signs.)
\end{theorem}

\begin{proof}
A non-residue class contains no squares, so by \eqref{eq:CN} and Lemma~\ref{lem:structure},
$0\le C_N=c_N(E)-\sum_\eta\beta_\eta\sqrt{x_\eta}B_N(x_\eta)$ for $N\equiv b\pmod q$. Multiply by $e^{-\pi Nv}$ and sum:
\[
0\le\sum_{N\equiv b\,(q)}c_N(E)e^{-\pi Nv}-\sum_\eta\beta_\eta\sqrt{x_\eta}\,T^{(x_\eta)}_{q,b}(v).
\]
The first sum vanishes if $E=\emptyset$; for $(q,b)=(3,2)$ it is $O(v^{-1/4})$ by Fact~\ref{lem:cusp-class2}. Multiply by
$v^{1/2}$ and let $v\to0^+$; by Proposition~\ref{prop:classsum}, $v^{1/2}T^{(x_\eta)}_{q,b}(v)\to\kappa_{q,b}(x_\eta)$, which is $0$ unless
$\eta\in\mathcal X_\QQ$. For the consequence, if $\lambda_{q,b}\ge0$ and $\sum\lambda_{q,b}\kappa_{q,b}(x_\eta)>0$ for $\eta\in\mathcal X_\QQ$, then
$0\ge\sum_\eta\beta_\eta\sqrt{x_\eta}\sum\lambda_{q,b}\kappa_{q,b}(x_\eta)$ is a sum of non-negative terms, so each vanishes.
\end{proof}

For $E\ne\emptyset$ and $q\ge5$ the argument needs the analogue of Fact~\ref{lem:cusp-class2} for the class $b\bmod q$, which
is not proved here (Section~\ref{sec:open}). For $E=\emptyset$, Fact~\ref{fact:zerosupport} gives more than part (a): every $\beta_\eta$
vanishes. Part (a) remains a statement about the class sums.

\begin{lemma}[A closed form on a dense family]\label{lem:D}
Let $d\equiv1\pmod{24}$ be prime, let $a\equiv5\pmod6$ be an integer with $d\nmid a$, and put $X:=a/d$, $x:=\sqrt X$. Then
\[
x\,\kappa_{3,2}(x)=\frac{(a/d)}{3\sqrt d},
\]
where $(a/d)$ is the Legendre symbol. The set of such $X$ with $(a/d)=+1$ is dense in $(0,\infty)$; hence
$\{X\in\QQ_{>0}:\kappa_{3,2}(\sqrt X)>0\}$ is dense.
\end{lemma}

The proof (Appendix~\ref{app:proofs-lemD}) evaluates the four values of $g$ in Proposition~\ref{prop:classsum} by quadratic
Gauss sums, and obtains the density from the P\'olya--Vinogradov inequality.

\begin{proposition}[Irrational extras are invisible at leading order]\label{prop:invisible}
If $x^2\notin\QQ$, then $T^{(x)}_{q,b}(v)=o(v^{-1/2})$ as $v\to0^+$, for every odd prime $q$ and every non-residue $b\bmod q$. Consequently
an extra atom $\eta$ with $x_\eta^2\notin\QQ$ makes no contribution to the inequalities of Theorem~\ref{thm:extra-primes}.
\end{proposition}

\begin{proof}
This is the irrational case of Proposition~\ref{prop:classsum}.
\end{proof}

The reading of this proposition is an obstruction: no argument at leading order in the class sums can exclude irrational-type
extras. Shifted prime powers $x^2=n^2\pm\delta$,
and the spurious points of near-optimal test functions observed numerically, are of this type. Heuristically, such extras are
visible at intermediate scales: near a good rational approximation of $x^2$ the class sum behaves like that of the rational
point (Numerical observation~\ref{obs:routeC}).

\begin{lemma}[Continuity of $G_x$ in $x$]\label{lem:Gx-cont}
For fixed $\tau\in\mathbb H$, the map $x\mapsto G_x(\tau)$ is continuous on $(0,\infty)$.
\end{lemma}

The proof is in Appendix~\ref{app:proofs-Gxcont}. It uses the contour-integral construction of the modular integral in
\cite[\S3.3]{BRS}, dominated convergence on a fundamental domain, and the transport of continuity by the cocycle relation.

\begin{proposition}[Lone extras sit on a nowhere-dense set]\label{prop:C0}
Let $\mathfrak B$ be the set of $X\in(4,\infty)\setminus\ZZ$ such that $T^{(\sqrt X)}_{q,b}(v)\le0$ for every $v>0$, every odd prime $q$ and every
non-residue $b\bmod q$.
\begin{enumerate}[label=\textup{(\alph*)}]
\item If $(\mu,\nu)\in\Kad$ with $\mu$ carried by $\Zz$ and $\nu$ carried by $\BRS\cup\{\eta\}$, where $\eta\notin\BRS$ and $\nu(\{\eta\})>0$, then
$x_\eta^2\in\mathfrak B$.
\item $\mathfrak B$ is closed relative to $(4,\infty)\setminus\ZZ$, and nowhere dense in $\RR$.
\end{enumerate}
By Fact~\ref{fact:zerosupport}, the hypothesis of \textup{(a)} never holds; \textup{(b)} is a statement about theta class sums.
\end{proposition}

\begin{proof}
(a) Here $x_\eta\ge2$ and $x_\eta^2\notin\NN$, so $X:=x_\eta^2\in(4,\infty)\setminus\ZZ$. With $E=\emptyset$ and $\mathcal X=\{\eta\}$,
\eqref{eq:CN} gives $0\le C_N=-\beta_\eta\sqrt{x_\eta}B_N(x_\eta)$ at every non-square $N$, so $B_N(x_\eta)\le0$ for all $N$ in every
non-residue class, and $T^{(x_\eta)}_{q,b}(v)\le0$. (b) For fixed $v$, $q$, $b$, the map $X\mapsto T^{(\sqrt X)}_{q,b}(v)$ is a finite
combination of values $G_{\sqrt X}(\tau)$, hence continuous by Lemma~\ref{lem:Gx-cont}; so $\mathfrak B$ is an intersection of relatively
closed sets. If $\kappa_{q,b}(\sqrt X)>0$ for some $(q,b)$, then $X\notin\mathfrak B$ by Proposition~\ref{prop:classsum}. If the closure of
$\mathfrak B$ contained an open interval $I$, then $I\subset[4,\infty)$ and, by relative closedness, $I\setminus\ZZ\subset\mathfrak B$; but by
Lemma~\ref{lem:D}, $I\setminus\ZZ$ contains rationals $X$ with $\kappa_{3,2}(\sqrt X)>0$.
\end{proof}

Only $E=\emptyset$ is covered. Whether $\mathfrak B=\emptyset$ is open (Section~\ref{sec:open}).

%======================================================================
\section{Positive spanning by the BRS functions}\label{sec:signpatterns}
%======================================================================

The argument behind Fact~\ref{thm:Gfin} gives more than its statement: a positive spanning property of the values of the BRS
functions $U_m$, $m\equiv2\pmod3$, on the zero side. It is recorded here because it comes from the same class sums; it gives
no information on extra prime atoms (Remark~\ref{rem:signpattern-test}).

\begin{proposition}[Positive spanning by the BRS functions]\label{prop:signpattern}
Let $E\subset(0,\infty)\setminus\Zz$ be finite.
\begin{enumerate}[label=\textup{(\alph*)}]
\item For every real vector $a=(a_e)_{e\in E}\ne0$ there is $m\equiv2\pmod3$ with $\sum_{e\in E}a_eU_m(e)>0$. Equivalently, the vectors
$(U_m(e))_{e\in E}$, $m\equiv2\pmod3$, positively span $\RR^E$, and finitely many of them suffice.
\item For every real vector $a=(a_e)_{e\in E\cup\{0\}}\ne0$ with $a_0\ge0$ there is $m\equiv2\pmod3$ with $\sum_{e\in E\cup\{0\}}a_eU_m(e)>0$.
Equivalently, the unit vector $\delta_0$ at the point $0$ lies in the interior of the convex cone in $\RR^{E\cup\{0\}}$ generated by the
vectors $(U_m(e))_{e\in E\cup\{0\}}$, $m\equiv2\pmod3$.
\end{enumerate}
\end{proposition}

\begin{proof}
Let $E'$ be $E$ in (a) and $E\cup\{0\}$ in (b), and suppose that $\sum_{e\in E'}a_eU_m(e)\le0$ for every $m\equiv2\pmod3$. Put
$s_e:=\frac14+\frac{ie}2$ and $\tilde\varrho_e:=a_e\Xi(e)/(\frac14+e^2)$ for $e\in E'$ (so $\tilde\varrho_0=4a_0\Xi(0)$), and
$c_N:=\sum_{e\in E'}\tilde\varrho_e\alpha_N(s_e)$. At $s=\frac12+ie$ one has $\frac12s(s-1)=-\frac12(\frac14+e^2)$, so
$\GR(s)\zeta(s)=-2\Xi(e)/(\frac14+e^2)$, and $\frac s2=s_e$. Hence \eqref{eq:Um} gives
\[
-2\pi m^{1/4}\sum_{e\in E'}a_eU_m(e)=\sum_{d^2\mid m}\boldsymbol\mu(d)\,c_{m/d^2}=:c^*_m,
\]
and $c^*_m\ge0$ for every $m\equiv2\pmod3$. By M\"obius inversion, $c_N=\sum_{k^2\mid N}c^*_{N/k^2}$. If $N\equiv2\pmod3$ and $k^2\mid N$, then
$3\nmid k$, so $N/k^2\equiv2\pmod3$; hence $c_N\ge0$ for every $N\equiv2\pmod3$. In (b), $\tilde\varrho_0\ge0$ because $\Xi(0)>0$.
Fact~\ref{lem:class2-landau}, applied with $E_+=E'$ and with the weight $2a_0$ in place of $a_0$ (which gives the same $\tilde\varrho_0$), yields
$\tilde\varrho_e=0$ for every $e\in E'$, and $\Xi(e)\ne0$ because $e\notin\Zz$; so $a=0$.

The equivalences are facts about the convex cone $K\subset\RR^n$ generated by the vectors in question and its dual cone
$K^*=\{a:\langle a,v\rangle\ge0\text{ for all }v\in K\}$. In (a) the statement says that $K^*=\{0\}$. Then $K$ is dense, hence equal to $\RR^n$,
because a convex set and its closure have the same interior; so each of the $2n$ vectors $\pm\delta_e$ is a finite non-negative
combination of generators. In (b) the statement, applied to $-a$, says that $a_0>0$ for every non-zero $a\in K^*$, and this holds if and only
if $\delta_0$ lies in the interior of $K$.
\end{proof}

\begin{remark}[Testing admissible pairs]\label{rem:signpattern-test}
Let $(\mu,\nu)$ be an admissible pair with $\mu$ carried by $\Zz\cup E$, $E$ finite and symmetric, and $\nu$ carried by $\BRS$
(Lemma~\ref{lem:structure} with $\mathcal X=\emptyset$), and put $a_e:=\mu(\{e\})\ge0$. Let $F=\sum_{m\equiv2\,(3)}x_mU_m$ be a finite
combination with $x_m\ge0$. Then $\Arch(F)=0$ by \eqref{eq:W-BRS}, since no $m\equiv2\pmod3$ is a square; $F$ vanishes at every $t_\rho$;
and $\hF$ equals $x_m$ at the node $\log m/4\pi$ and $0$ at the other nodes. So the explicit formula for the pair gives
$\sum_{e\in E}a_eF(e)=-\frac1\pi\sum_mx_m\,\nu(\{\frac{\log m}{4\pi}\})\le0$. If $0\in E$, Proposition~\ref{prop:signpattern}(b), applied to
$E_+\setminus\{0\}$, gives such an $F$ with $F(0)>0$ and $F=0$ on $E\setminus\{0\}$; so $a_0F(0)\le0$ and $a_0=0$. Then, for each
$e_0\in E_+\setminus\{0\}$, part (a), applied to $E_+\setminus\{0\}$, gives such an $F$ with $F(e_0)>0$ and $F=0$ on $E\setminus\{0,\pm e_0\}$;
so $2a_{e_0}F(e_0)\le0$ and $a_{e_0}=0$. Hence $\mu$ is carried by $\Zz$, and $(\mu,\nu)=\pz$ by Lemma~\ref{lem:structure}: for admissible
pairs, testing with explicit functions recovers Fact~\ref{thm:Gfin}.

These test functions change sign on $\RR$: $\int F=\hF(0)=0$, because the node $0=\log1/4\pi$ belongs to the square $m'=1$ and
$\widehat{U_m}(0)=0$ for $m\ge2$. So they are not in the cone of non-negative functions with non-negative transform beyond the gap. In the
terms of Lemma~\ref{lem:P0} they are node-positive, but nothing controls the sign of $\hF$ at a point $\eta\notin\BRS$, so they give no
information on extra prime atoms.
\end{remark}

%======================================================================
\section{Countably many extra zeros: an outline}\label{sec:countable}
%======================================================================

\begin{remark}[Countably many extra zeros]\label{rem:Gsum}
We had proposed to extend Fact~\ref{thm:Gfin} to countable, locally finite sets $E$ with
$\sum_{e\in E}\mu(\{e\})\,|\zeta(\frac12+ie)|\,(1+|e|)^{-1/4+\delta}<\infty$ for some $\delta>0$, which would need control, uniform in $e$,
of the cusp expansions at the regular points $\frac23$ and $\frac12$. This is now proved in \cite[Theorem~3.8]{PaperI}, by a theorem of Astra, under the
weaker hypothesis $\sum_{e\in E}\mu(\{e\})\,|\zeta(\frac12+ie)|\,(1+|e|)^{-1/2}<\infty$ (more generally, when the partial sums of
$\mu(\{e\})|\zeta(\frac12+ie)|$ over $0\le e\le T$ are $o(\sqrt T)$) and without local finiteness: the intermediate cusp terms are bounded
in aggregate, and a mean-value argument replaces Landau's theorem. Without such a summability condition, for instance for a positive
density of extra zeros with masses of order $1$, the question remains open.
\end{remark}

%======================================================================
\section{Numerical observations}\label{sec:observations}
%======================================================================

The two observations below were computed but not certified; they are not used in any proof.

\begin{observation}[Rational extras]\label{obs:rational-extras}
For each of the $126$ rationals $X=a/d\in(\frac12,12)\setminus\ZZ$ with $2\le d\le6$, some class $(q,b)$ with $q\in\{3,5,7,11,13\}$ has
computed value $x\,\kappa_{q,b}(x)>0$; the smallest such value is $0.004278$, at $X=7/6$, $(q,b)=(13,6)$. For the $7875$ pairs of these
values, a non-negative combination of the $17$ class vectors with $q\le13$ is positive at both points for $7872$ pairs, and for the
remaining $3$ once the classes with $q\in\{17,19,23\}$ are added. The values for $q=3$ are exact Gauss-sum evaluations; for $q\ge5$ they
were obtained by numerical extrapolation of $v^{1/2}T^{(x)}_{q,b}(v)$ as $v\downarrow3\cdot10^{-6}$, with no closed form and no interval
certification. Only $X>4$ corresponds to positions in $[\xin2,\infty)$. If certified, Theorem~\ref{thm:extra-primes}(a) would exclude a lone
rational extra atom at each such $X$ when $E=\emptyset$; for admissible pairs that case is now settled by Fact~\ref{fact:zerosupport}, so
these values concern the class sums only.
\end{observation}

\begin{observation}[Irrational extras at intermediate scales]\label{obs:routeC}
Let $E=\emptyset$. For each of $12$ irrational-type positions (seven shifted prime powers $x^2=n^2\pm\delta$ with $n\in\{3,5,7,8,11\}$ and
$\delta\in[0.003,0.067]$, and $x^2\in\{4+(\sqrt2-1),\,5+(\phi-1),\,9+1/e,\,16+\sqrt3/10,\,30+\sqrt7/10\}$, $\phi$ the golden ratio), the computed
function $v^{1/4}T^{(x)}_{3,2}(v)$ reaches a positive value, between $+0.10$ and $+0.81$, somewhere in $10^{-9}\le v\le10^{-2.75}$; and for all
$15$ pairs formed from six of these positions the corresponding linear feasibility problem has positive margin ($0.08$--$0.67$). The class
sums were computed from the exact cocycle formula with exponentially small terms dropped and $\vartheta$ evaluated to $25$ digits; they are
not interval-certified. If certified, Lemma~\ref{lem:structure} would exclude a lone extra atom (resp.\ these pairs) at these positions;
with $E=\emptyset$ this is now implied by Fact~\ref{fact:zerosupport}, so the computation concerns the class sums only.
\end{observation}

%======================================================================
\section{Open problems}\label{sec:open}\label{ssec:open}
%======================================================================

\begin{enumerate}[label=\textup{(\arabic*)}]
\item \emph{Theta class sums.} Is the set $\mathfrak B$ of Proposition~\ref{prop:C0} empty? For admissible pairs without extra zeros the
question is settled by Fact~\ref{fact:zerosupport}, but it remains a question about theta class sums: Theorem~\ref{thm:extra-primes} and
Proposition~\ref{prop:C0} leave a closed, nowhere-dense set of positions, and irrational positions are invisible to the leading-order
argument (Proposition~\ref{prop:invisible}). This seems to need theta-sum asymptotics along continued fractions and a
Duffin--Schaeffer-type step \cite{DS41,KoMa20}.
\item \emph{Extra zeros and extra prime atoms together.} Exclude extra prime mass when $\mu$ has finitely many extra zeros and $\nu$ is
arbitrary. Lemma~\ref{lem:P0} gives one route: a single node-positive test function that vanishes on $E$ and is positive at the
position of the atom excludes it.
\item \emph{Extra zeros and other classes.} For $E\ne\emptyset$ and $q\ge5$, the analogue of Fact~\ref{lem:cusp-class2} for the
non-residue classes $b\bmod q$, which would extend Theorem~\ref{thm:extra-primes}(b) to these classes.
\item \emph{Infinitely many extra zeros} without the summability condition of Section~\ref{sec:countable}.
\item \emph{Certification} of the computations of Numerical observations~\ref{obs:rational-extras} and~\ref{obs:routeC}.
\end{enumerate}

\section*{AI-use statement}
This note was produced with substantial use of AI. The mathematical arguments, the numerical computations and the text were generated by instances of Claude Opus 5.5
(Anthropic) working as coordinated agents: one set of agents developed the arguments, separate agents acted as
independent hostile referees of each component, and the referees' requested repairs were incorporated. The zero-side
support theorem and its extension to countably many extra zeros (Theorems~3.6 and~3.8 of the companion paper), which this note
cites, were contributed by Astra (OpenAI) and checked by independent referee agents before inclusion. The
author initiated and directed this project, designed its verification process, and reviewed the results and the
verification record, but did not independently check every step of the written proofs. The author takes full
responsibility for the content of this note, including any errors, and welcomes independent verification and
corrections.

\appendix
%======================================================================
\section{Proofs}\label{app:proofs}
%======================================================================

\subsection{Structure of pairs with extra atoms}\label{app:proofs-structure}
\begin{proof}[Proof of Lemma~\textup{\ref{lem:structure}}]
Since $U_m$ vanishes on $\Zz$, $\widehat{U_m}(\frac{\log m'}{4\pi})=\delta_{mm'}$ and \eqref{eq:W-BRS} holds, testing on $U_m$ gives
\[
\sum_{e\in E}a_eU_m(e)+\frac1\pi\Big(\nu\big(\{\tfrac{\log m}{4\pi}\}\big)+\sum_\eta\beta_\eta\widehat{U_m}(\eta)\Big)=\frac1\pi\Lambda(\sqrt m)m^{-1/4}.
\]
Multiply by $\pi m^{1/4}$. By \eqref{eq:Um} and $\GR(\frac12+it)\zeta(\frac12+it)=-2\Xi(t)/(\frac14+t^2)$,
$\pi m^{1/4}a_eU_m(e)=-\frac{a_e\Xi(e)}{2(\frac14+e^2)}\sum_{d^2\mid m}\boldsymbol\mu(d)\alpha_{m/d^2}(s_e)$. The atoms at $\pm e$ give equal
terms, because $a_{-e}=a_e$, $\Xi$ is even and $\alpha_N(s_{-e})=\alpha_N(\overline{s_e})=\alpha_N(s_e)$; so
$\pi m^{1/4}\sum_ea_eU_m(e)=-\sum_{d^2\mid m}\boldsymbol\mu(d)c_{m/d^2}(E)$. With $m^{1/4}\widehat{U_m}(\eta)=\sqrt{x_\eta}u_m(x_\eta)$ we get
\[
\ell_m=\Lambda(\sqrt m)\one_\square(m)+\sum_{d^2\mid m}\boldsymbol\mu(d)\Big(c_{m/d^2}(E)-\sum_\eta\beta_\eta\sqrt{x_\eta}B_{m/d^2}(x_\eta)\Big).
\]
M\"obius inversion over squares ($f_m=\sum_{d^2\mid m}\boldsymbol\mu(d)g_{m/d^2}$ for all $m$ if and only if
$g_N=\sum_{k^2\mid N}f_{N/k^2}$ for all $N$) and $\sum_{k^2\mid N}\Lambda(\sqrt{N/k^2})\one_\square(N/k^2)=\sum_{k\mid M}\Lambda(M/k)=\log M$ for
$N=M^2$ give \eqref{eq:CN}. Now let $\mathcal X=\emptyset$. If $\nu$ lives on $[\xin2,\infty)$ then $\ell_1=\ell_2=\ell_3=0$, and
\eqref{eq:CN} at $N\le3$ gives $c_N=0$. The remaining statements for $\mathcal X=\emptyset$ come from testing on $V_{\rho,j}$,
whose transform vanishes on $\BRS$, with \eqref{eq:W-BRS}, and from the uniqueness of $\zeta$'s pair on its own support; they
are proved in \cite{PaperI}. Finally, if $(\mu,\nu)\in\Kad$ then $\ell_m\ge0$, so $C_N\ge0$.
\end{proof}

\subsection{The theta function near rational points}\label{app:proofs-theta}
\begin{proof}[Proof of Lemma~\textup{\ref{lem:theta-local}}]
(a) Since $e^{\pi in^2r}$ depends only on $n\bmod 2c$, Poisson summation in each residue class gives
\[
\vartheta(r+z)=\sum_{\nu\bmod 2c}e^{\pi i\nu^2r}\sum_{m\in\ZZ}e^{\pi i(\nu+2cm)^2z}
=\frac{(z/i)^{-1/2}}{2c}\sum_{\nu\bmod2c}e^{\pi i\nu^2r}\sum_{k\in\ZZ}e^{\pi ik\nu/c}e^{-\pi ik^2/(4c^2z)} .
\]
The term $k=0$ is $(z/i)^{-1/2}(2c)^{-1}\sum_{\nu\bmod2c}e^{\pi i\nu^2a/c}=g(r)(z/i)^{-1/2}$: if $ac$ is even, $\nu\mapsto e^{\pi i\nu^2a/c}$ has period
$c$; if $ac$ is odd, the terms $\nu$ and $\nu+c$ cancel. The terms $k\ne0$ have modulus at most $e^{-\pi k^2Y/(4c^2)}$. In a sector
$|\Re z|\le M\Im z$ one has $Y=\Im z/|z|^2\ge(1+M^2)^{-1/2}|z|^{-1}$.
(b) First, $|\vartheta(r+\delta+i\eps)-\vartheta(r+i\eps)|\le\sum_n\pi n^2|\delta|e^{-\pi n^2\eps}\ll M\eps^2\eps^{-3/2}$. Next, by Weyl's
equidistribution theorem for $n^2r/2$ \cite[Ch.~1, Theorems~2.1 and~3.2]{KN74}, $S(n):=\sum_{|k|\le n}e^{\pi ik^2r}=o(n)$. Partial summation gives
$\vartheta(r+i\eps)=\sum_{n\ge0}S(n)\big(e^{-\pi n^2\eps}-e^{-\pi(n+1)^2\eps}\big)$; given $\eta>0$, split at $n_0$ with $|S(n)|\le\eta n$ for $n\ge n_0$ and use
$\sum_nn(e^{-\pi n^2\eps}-e^{-\pi(n+1)^2\eps})=\sum_{n\ge1}e^{-\pi n^2\eps}\le(2\eps)^{-1/2}$.
(c) Let $\tau=r+i\eps$, $\eps\to0^+$. Then $-1/\tau=-1/r+z'$ with $z'/i=\eps r^{-2}(1+O(\eps))$, so by (a) both sides of
$\vartheta(-1/\tau)=(\tau/i)^{1/2}\vartheta(\tau)$ are $\eps^{-1/2}$ times a constant plus $o(\eps^{-1/2})$: the constant is $g(-1/r)|r|$ on the left
and $(r/i)^{1/2}g(r)$ on the right. As $(r/i)^{1/2}=|r|^{1/2}e^{-i\pi\sgn(r)/4}$, the claim follows.
\end{proof}

\subsection{Class sums of an extra atom}\label{app:proofs-classsum}
\begin{proof}[Proof of Proposition~\textup{\ref{prop:classsum}}]
Fix $r=2j/q$ with $0\le j<q$. Reduce $r$ to $\infty$ by the following adaptation to cusps of the reduction algorithm of
\cite[\S3.1]{BRS} (which is stated for points of $\mathbb H$): translate by an even integer into
$[-1,1]$, then apply $S$, and repeat. On reduced fractions $a/c$ the translations preserve the parities of $a$ and $c$ and the
denominator, and $S:a/c\mapsto-c/a$ exchanges the types (even, odd) and (odd, even) and lowers the denominator, since $|a|<c$ in
$[-1,1]$ (the value $\pm1$, of type (odd, odd), does not occur). So the procedure ends at $0/1$ followed by a final $S$; it gives a
word $\gamma_r=g_m\cdots g_1$ with $g_m=S$, $\gamma_r r=\infty$, and $w_{m-1}r=0$. As $r=2j/q$ and $0/1$ have type (even, odd), the number
of $S$'s in $w_{m-1}$ is even, so $\chi(\gamma_r)=-1$ and $\chi(w_{m-1})=+1$. Every other $k$ with $g_k=S$ has $\rho_k:=w_{k-1}r\in\QQ\setminus\{0\}$.
Let $\tau=r+iv$. Each $w_{k-1}$ is a M\"obius map, finite at $r$, so $w_{k-1}\tau=\rho_k+i\lambda_kv+O(v^2)$ with
$\lambda_k:=w_{k-1}'(r)>0$; and $\Im\gamma_r\tau=1/(c_r^2v)$ where $c_r\ne0$ is the lower-left entry of $\gamma_r$. Apply \eqref{eq:unroll} to $G_x$.

\emph{The term $G_x|\gamma_r$.} It equals $-\vartheta(\tau)G_x(\gamma_r\tau)/\vartheta(\gamma_r\tau)=-B_0(x)\vartheta(\tau)+O(v^{-1/2}e^{-\pi/(c_r^2v)})$.

\emph{The term $k=m$.} Here $w:=w_{m-1}\tau\to0$ non-tangentially. Inverting both theta functions in $\mathfrak p_x(w)$,
$\vartheta(x^2w)=x^{-1}(w/i)^{-1/2}\vartheta(-1/(x^2w))$ and $x^{-1}\vartheta(w/x^2)=(w/i)^{-1/2}\vartheta(-x^2/w)$, gives
$\mathfrak p_x(w)=\frac1{2x}(w/i)^{-1/2}-\frac12+O(|w|^{-1/2}e^{-c/|w|})$. Since $(w/i)^{-1/2}/\vartheta(w)=1/\vartheta(-1/w)=1/\vartheta(\gamma_r\tau)$,
\[
(\mathfrak p_x|w_{m-1})(\tau)=\frac{\vartheta(\tau)}{\vartheta(w)}\,\mathfrak p_x(w)=\frac1{2x}\vartheta(\tau)-\frac12\,\frac{\vartheta(\tau)}{\vartheta(w_{m-1}\tau)}+O(v^{-1/2}e^{-c/v}),
\]
and the middle term is bounded: $\vartheta(\tau)/\vartheta(w_{m-1}\tau)$ is a product of factors $(w_{k-1}\tau/i)^{-1/2}$ over the $S$'s in $w_{m-1}$,
evaluated at points tending to $\rho_k\ne0$.

\emph{The terms $k<m$ with $g_k=S$.} Here $(\mathfrak p_x|w_{k-1})(\tau)=\chi(w_{k-1})j_k(\tau)\mathfrak p_x(w_{k-1}\tau)$, where
$j_k(\tau):=\vartheta(\tau)/\vartheta(w_{k-1}\tau)$ is analytic near $r$ with a value $j_k\ne0$ there. In $\mathfrak p_x(w_{k-1}\tau)$ the terms $-1$
and $(w_{k-1}\tau/i)^{-1/2}$ are bounded, while $\vartheta(x^2w_{k-1}\tau)$ and $\vartheta(w_{k-1}\tau/x^2)$ are theta values near the real points
$x^2\rho_k$ and $\rho_k/x^2$, approached in the manner of Lemma~\ref{lem:theta-local}(a),(b). If $x^2\notin\QQ$ these points are irrational and
the terms are $o(v^{-1/2})$. If $x^2\in\QQ$, Lemma~\ref{lem:theta-local}(a) gives
$\vartheta(x^2w_{k-1}\tau)=g(x^2\rho_k)(x^2\lambda_kv)^{-1/2}+O(v^{1/2})$ and $\vartheta(w_{k-1}\tau/x^2)=g(\rho_k/x^2)(\lambda_kv/x^2)^{-1/2}+O(v^{1/2})$.

\emph{Summation over the cusps.} The coefficient of $\vartheta(\tau)$ is $-(B_0(x)-\frac1{2x})$ at every cusp $2j/q$, and
$\sum_{j\bmod q}e^{-2\pi ijb/q}\vartheta(\frac{2j}q+iv)=q\sum_{n^2\equiv b\,(q)}e^{-\pi n^2v}=0$ since $b$ is a non-residue. All other terms are $O(1)$,
exponentially small, or theta values at the sampled points. Hence $T^{(x)}_{q,b}(v)=\kappa_{q,b}(x)v^{-1/2}+o(v^{-1/2})$ with $\kappa_{q,b}(x)=0$ if
$x^2\notin\QQ$, and with an error $O(1)$ and
\begin{equation}\label{eq:kappa-general}
\kappa_{q,b}(x)=\frac1q\sum_{j\bmod q}e^{-2\pi ijb/q}\sum_{k<m(j),\ g_k=S}\chi(w_{k-1})\,j_k\,\frac{\lambda_k^{-1/2}}2\Big[\frac{g(x^2\rho_k)}x+g(\rho_k/x^2)\Big]
\end{equation}
if $x^2\in\QQ$ (all quantities depending on $j$).

\emph{The case $(q,b)=(3,2)$.} For $j=0$ the word is $S$ and there is no term with $k<m$. For $j=1$ the word is $ST^2ST^2S$
(Section~\ref{ssec:cusp}) and the terms $k<m$ are $k=1$ ($w_0=1$, $\rho=\frac23$, $\chi=1$, $j_1=1$, $\lambda=1$) and $k=3$ ($w_2=T^2S$,
$\rho=\frac12$, $\chi=-1$, $j_3=(\frac2{3i})^{-1/2}=(\frac32)^{1/2}e^{i\pi/4}$, $\lambda=(\frac32)^2$). With $X=x^2$ their contribution is
\[
K:=\frac12\Big[\frac{g(\frac{2X}3)}x+g\big(\tfrac2{3X}\big)\Big]-\frac{e^{i\pi/4}}{\sqrt6}\Big[\frac{g(\frac X2)}x+g\big(\tfrac1{2X}\big)\Big].
\]
By Lemma~\ref{lem:theta-local}(c), $g(\frac2{3X})=(\frac{3X}2)^{-1/2}e^{i\pi/4}g(-\frac{3X}2)$ and $g(\frac1{2X})=(2X)^{-1/2}e^{i\pi/4}g(-2X)$, which turns $K$
into the expression in Proposition~\ref{prop:classsum}. For $j=2$, $\frac43\equiv-\frac23\pmod2$ and the $B_N(x)$ are real, so
$G_x(\frac43+iv)=\overline{G_x(\frac23+iv)}$ and the contribution is $\overline K$. With $e^{-4\pi i/3}=\omega$ and $e^{-8\pi i/3}=\bar\omega$ we get
$\kappa_{3,2}(x)=\frac13(\omega K+\bar\omega\overline K)=\frac23\Re(\omega K)$.

\emph{Integers.} If $X=m\in\NN$, then $B_N(\sqrt m)=\sum_{k\ge1}b_N(k\sqrt m)=\#\{k\ge1:k^2m=N\}$, so $G_x=\frac12(\vartheta(m\tau)-1)$ and
$T^{(x)}_{3,2}(v)=\sum_{k\ge1,\,k^2m\equiv2\,(3)}e^{-\pi k^2mv}$, which is $\frac13x^{-1}v^{-1/2}+O(1)$ if $m\equiv2\pmod3$ and $0$ otherwise. Evaluating
$K$ directly (for instance $m=5$: $g(\frac{10}3)=-i/\sqrt3$, $g(-\frac{15}2)=g(\frac52)=\frac{1+i}2$, $g(-10)=1$, so $K=-i/(\sqrt3x)$) gives the same values.
\end{proof}

\subsection{A closed form on a dense family}\label{app:proofs-lemD}
\begin{proof}[Proof of Lemma~\textup{\ref{lem:D}}]
Let $G(c;n):=\sum_{k\bmod n}e^{2\pi ick^2/n}$. We use the standard facts on quadratic Gauss sums \cite[\S3.4]{IK04}: $G(c;n)=(c/n)G(1;n)$ for odd $n$ and $(c,n)=1$, with $G(1;n)=\sqrt n$ or $i\sqrt n$ according as
$n\equiv1$ or $3\pmod4$; $G(c;n_1n_2)=G(cn_2;n_1)G(cn_1;n_2)$ for coprime $n_1,n_2$; and $G(c;4)=2(1+i^c)$ for odd $c$. Since $a$ is odd,
$3\nmid a$, $d\nmid a$ and $d\ne3$, the four points $\frac{2a}{3d}$, $-\frac{2a}d$, $\frac a{2d}$, $-\frac{3a}{2d}$ are in lowest terms with
even numerator or even denominator. Summing over $n\bmod 2c$ when the denominator $c$ is even, one finds
\[
g\big(\tfrac{2X}3\big)=\frac{G(a;3d)}{3d},\quad g(-2X)=\frac{G(-a;d)}d,\quad g\big(\tfrac X2\big)=\frac{G(a;4d)}{4d},\quad g\big(-\tfrac{3X}2\big)=\frac{G(-3a;4d)}{4d}.
\]
Now $d\equiv1\pmod{12}$ gives $(-1/d)=(3/d)=(d/3)=1$, and $a\equiv2\pmod3$ gives $(a/3)=-1$. Hence
$G(a;3d)=G(ad;3)G(3a;d)=-i\sqrt3\,(a/d)\sqrt d$, $G(-a;d)=(a/d)\sqrt d$, and, as $i^{ad}=i^{-3ad}=i^a$,
$G(a;4d)=G(4a;d)G(ad;4)=2(1+i^a)(a/d)\sqrt d=G(-12a;d)G(-3ad;4)=G(-3a;4d)$. So $g(-\frac{3X}2)=g(\frac X2)$,
$g(\frac{2X}3)=-i(a/d)/\sqrt{3d}$, $g(-2X)=(a/d)/\sqrt d$, and $K=-i(a/d)/(x\sqrt{3d})$. Then $\Re(\omega K)=\frac{\sqrt3}2\cdot\frac{(a/d)}{x\sqrt{3d}}$,
which gives the formula. For density, fix an interval $(\alpha,\alpha')\subset(0,\infty)$. For a prime $d\equiv1\pmod{24}$ there are
$(\alpha'-\alpha)d/6+O(1)$ integers $a\equiv5\pmod6$ with $a/d\in(\alpha,\alpha')$, of which $O(1)$ are divisible by $d$; writing
$a=6a'+5$, the sum of $(a/d)$ over them is a sum of a non-principal character over an interval of consecutive integers $a'+5\cdot\bar6$,
hence $O(\sqrt d\log d)$ by the P\'olya--Vinogradov inequality \cite[Ch.~12]{IK04}. By Dirichlet's theorem there are arbitrarily
large such primes $d$, and for large $d$ some of these $a$ have $(a/d)=+1$.
\end{proof}

\subsection{Continuity in \texorpdfstring{$x$}{x}}\label{app:proofs-Gxcont}
\begin{proof}[Proof of Lemma~\textup{\ref{lem:Gx-cont}}]
Let $\phi_x(z)=\frac12(\vartheta(x^2z)-1)$ and fix $0<a<b$. For $x\in[a,b]$ and $z\in\mathbb H$,
$|\phi_x(z)|\le\sum_{n\ge1}e^{-\pi n^2a^2\Im z}\le(2a)^{-1}(\Im z)^{-1/2}$, and $x\mapsto\phi_x(z)$ is continuous.

By \cite[Proposition~3.3]{BRS}, the function $G_x=F^-_{1/2}(\cdot,\phi_x)$ is the analytic continuation to $\mathbb H$ of the
integral $\frac12\int_{\ell_0}\mathcal K(\tau,z)\phi_x(z)\dd z$ $(\tau\in\mathfrak F)$, where $\mathfrak F=\{|\Re\tau|<1,\ |\tau|>1\}$, $\ell_0$ is the upper unit semicircle
from $-1$ to $1$, and $\mathcal K=\mathcal K^-_{1/2}$ is the modular kernel of \cite[\S3.3]{BRS}. For fixed $\tau$, $z\mapsto\mathcal K(\tau,z)$ is
meromorphic in $\mathbb H$ with poles only on $\Gamma_\vartheta\tau$, and it decays like $e^{-c/\Im z}$ as $z\to\pm1$ non-tangentially. The continuation
is described in the proof of that proposition: the same integral defines it across the vertical sides $\Re\tau=\pm1$ of $\mathfrak F$ (no
point of $\Gamma_\vartheta\tau$ lies on $\ell_0$ for such $\tau$); and near $\ell_0$, $G_x=\frac12\int_{\ell_1}\mathcal K(\cdot,z)\phi_x(z)\dd z+\phi_x$ on the region
between a path $\ell_1$ from $-1$ to $1$ above $\ell_0$, meeting the real axis non-tangentially, and its image under $S$. Hence every
point $w$ of $\overline{\mathfrak F}\cap\mathbb H$ has a neighbourhood on which $G_x(w)$ is given by one of these integrals (plus $\phi_x(w)$), with
an integrand dominated, for all $x\in[a,b]$, by $|\mathcal K(w,z)|(2a)^{-1}(\Im z)^{-1/2}$, which is integrable along the path. By dominated
convergence, $x\mapsto G_x(w)$ is continuous on $[a,b]$ for every such $w$.

For general $\tau\in\mathbb H$ choose $\gamma\in\Gamma_\vartheta$ with $\gamma\tau\in\overline{\mathfrak F}$ \cite[\S3.1]{BRS}. By \eqref{eq:unroll},
$G_x(\tau)=\chi(\gamma)\frac{\vartheta(\tau)}{\vartheta(\gamma\tau)}G_x(\gamma\tau)+\sum_{k:\,g_k=S}(\mathfrak p_x|w_{k-1})(\tau)$. The first term is continuous in $x$ by
the previous paragraph, and the finitely many others are values of $\mathfrak p_x$, which is continuous in $x$, at points independent of $x$,
times factors independent of $x$.
\end{proof}

\bibliographystyle{amsplain-doi}
\bibliography{refs}

\end{document}
